\subsection{Transforming a Graph: Alternate Method} 
An alternate method of transforming a graph involves looking at the \makeblank[1in] and \makeblank[1in] of the original function. You might find this method helpful if you sometimes lose track.
\begin{Example}
Below is the graph of $y=f(x)$. Sketch the graph of $y=\frac{1}{2} f\left(-\frac{1}{3}x+2\right)-3$
\begin{icpic}[x=.5\linespace,y=.5\linespace]
\useasboundingbox (-7,-5) rectangle (19,7);
\setgraph{7}{19}{5}{7}
\GraphSmall
\draw[graphend] (-4,2) -- (-1,6) -- (3,4) -- (4,-2);
\draw[graphend] (3,4) node[anchor=west]{$y=f(x)$};
\end{icpic}
\begin{cpic}{}{6}
	\draw (-5,-2) -- (5,-2);
	\foreach \x in {-3,0,3}
		\draw (\x,-1) -- (\x,-6);
	\foreach \x/\lbl in {-4/$x$,-1.5/input,1.5/output,4/$y$}
		\regtext{\x}{-2}{\lbl};
\end{cpic}
\end{Example}
